\section{Synthèse du système sur FPGA (partie optionnelle)}
\label{sec:fpga}

Le modèle de la partie~\ref{sec:systeme} n'est que simulable. Le sujet demande les étapes
nécessaires pour le rendre synthétisable sur un FPGA Xilinx et, si le temps le permet, de les
mettre en œuvre. Nous ne disposions ni de carte Xilinx ni de Vivado, mais d'une carte équipée
d'un FPGA Intel, présentée à la section~\ref{sec:environnement}. Les étapes sont les mêmes
pour les deux constructeurs ; seuls les outils et le nom des ressources changent
(tableau~\ref{tab:xilinx-intel}). Les fichiers sont dans \texttt{exo/p6}.

\begin{table}[H]
\centering
\begin{tabular}{lll}
\toprule
 & Xilinx & Intel (carte utilisée) \\
\midrule
Outil & Vivado & Quartus II 13.1 \\
Contraintes & fichier \texttt{.xdc} & fichiers \texttt{.qsf} (broches) et \texttt{.sdc} (horloges) \\
Mémoire interne & blocs BRAM de 36~Kbit & blocs M9K de 9~Kbit \\
Multiplieurs & blocs DSP48 & multiplieurs $18\times18$ \\
Fichier de configuration & \texttt{.bit} & \texttt{.sof} \\
\bottomrule
\end{tabular}
\caption{Correspondance entre les flots Xilinx et Intel}
\label{tab:xilinx-intel}
\end{table}

\subsection{Environnement de travail et carte de développement}
\label{sec:environnement}

Le travail a été fait sur un PC sous Windows~11, avec les logiciels du
tableau~\ref{tab:environnement}. La carte utilisée est la carte de développement
\og New Starter\fg{} d'ALIENTEK (figure~\ref{fig:p6-carte-vue}). Elle est construite autour
d'un FPGA Intel Cyclone IV E EP4CE10F17C8 : 10\,320 éléments logiques, 46 blocs mémoire M9K
(414~Kbit au total), 23 multiplieurs $18\times18$ et 2 PLL, dans un boîtier BGA de
256 broches, en classe de vitesse 8. La carte comporte aussi une SDRAM, un port série et des
interfaces HDMI et LCD, que nous n'avons pas utilisés : le projet n'emploie que l'oscillateur
de 50~MHz, le bouton RESET, les quatre touches KEY0 à KEY3, les quatre LED et les six
afficheurs sept segments. Les numéros de broches sont pris dans le tableau de broches fourni
avec la carte. Le FPGA est programmé par JTAG, avec un câble USB-Blaster relié au PC.

\begin{table}[H]
\centering
\begin{tabular}{lll}
\toprule
Élément & Version ou référence & Utilisation \\
\midrule
Système & Windows 11 (10.0.26200) & \\
Quartus II & 13.1.0 Build 162, 64 bits & synthèse, placement-routage, timing, programmation \\
ModelSim & SE-64 2020.4 & simulation \\
Carte & ALIENTEK New Starter & FPGA Cyclone IV E EP4CE10F17C8 \\
Câble & USB-Blaster & programmation par JTAG \\
\bottomrule
\end{tabular}
\caption{Environnement de travail}
\label{tab:environnement}
\end{table}

\capture[0.8\linewidth]{p6-carte-vue.jpg}{Photo de la carte New Starter reliée au PC par le
câble USB-Blaster}
{La carte ALIENTEK New Starter (Cyclone IV E EP4CE10F17C8) et son câble USB-Blaster}{fig:p6-carte-vue}

\subsection{Étapes nécessaires}

\begin{enumerate}
\item \textbf{Repérer les constructions non synthétisables.} Les blocs VHDL des parties 2 et 3
le sont déjà : la RAM à adresse de lecture mémorisée correspond à une mémoire interne
synchrone. Le processeur, en revanche, attend plusieurs fronts dans un même bloc
\texttt{always} et utilise \texttt{\$display}, \texttt{\$finish}, une chaîne \texttt{string}
et l'option \texttt{+prog} : ces éléments n'ont pas d'équivalent matériel.
\item \textbf{Réécrire le processeur en machine à états synchrone} : un seul bloc
\texttt{always} sur le front montant, une instruction découpée en états.
\item \textbf{Fixer la mémoire programme} : une ROM initialisée à la synthèse par
\texttt{\$readmemh} avec un nom de fichier constant.
\item \textbf{Traiter le bus trois états interne} : un FPGA n'a pas de tampons trois états à
l'intérieur de la puce, seulement sur ses broches. L'outil remplace les quatre sorties de
\texttt{quadram} par un multiplexeur ; le code VHDL peut rester inchangé.
\item \textbf{Choisir l'horloge et écrire les contraintes} : fréquence compatible avec le
chemin le plus long, contraintes d'horloge pour l'analyse temporelle.
\item \textbf{Écrire un niveau supérieur propre à la carte} : bouton de reset resynchronisé,
affichage du résultat, affectation des broches d'après la documentation de la carte.
\item \textbf{Synthétiser, placer-router, vérifier le timing}, générer le fichier de
configuration, programmer le FPGA par JTAG et tester sur la carte.
\end{enumerate}

\subsection{Processeur synthétisable}

Le nouveau processeur, \texttt{riscvproc\_fpga.v}, garde les instructions et le bus de la
partie~\ref{sec:systeme}. Il est décrit par une machine à quatre états
(code~\ref{lst:fsm}) :
\begin{itemize}
\item \texttt{FETCH} : la ROM lit l'instruction à l'adresse \texttt{pc} ; la lecture étant
synchrone, l'instruction est disponible à la période suivante ;
\item \texttt{EXEC} : l'instruction est exécutée ; pour \texttt{sw}, l'adresse, la donnée et
\texttt{we} sont présents sur le bus et \texttt{quadram} écrit au front montant qui termine
la période ;
\item \texttt{MEM} : uniquement pour \texttt{lw} ; l'adresse a été mémorisée par
\texttt{quadram} à la fin de \texttt{EXEC}, la donnée lue est copiée dans le registre ;
\item \texttt{HALT} : atteint par \texttt{ebreak}, qui remplace \texttt{\$finish}.
\end{itemize}
Une instruction de calcul ou \texttt{sw} dure donc deux périodes, \texttt{lw} trois. Pendant
\texttt{EXEC} et \texttt{MEM}, l'adresse est calculée à partir de la même instruction et des
mêmes registres : elle reste stable, ce qui garantit que le décodeur de \texttt{quadram}
sélectionne la bonne RAM. Une instruction inconnue est ignorée et met à 1 la sortie
\texttt{illegal}, à la place du message d'erreur.

\begin{lstlisting}[style=verilog, caption={Machine à états du processeur synthétisable (extrait)}, label={lst:fsm}]
always @(posedge clock) begin
  if (reset) begin
    ...                                    // pc, registers and state cleared
  end else begin
    case (state)
      FETCH: state <= EXEC;
      EXEC: begin
        state <= FETCH;
        pc    <= pc + 12'd1;
        case (opcode)
          ARITHMETIC:
            if (!r_valid)     illegal  <= 1'b1;
            else if (rd != 0) regs[rd] <= r_result;
          LOAD:                            // lw: wait for the quadram
            if (f3 == 3'b010) begin state <= MEM;  pc <= pc; end
            else              illegal <= 1'b1;
          EBREAK:             begin state <= HALT; pc <= pc; end
          ...
        endcase
      end
      MEM: begin                           // lw: the word read is valid now
        if (rd != 0) regs[rd] <= dataToRead;
        pc    <= pc + 12'd1;
        state <= FETCH;
      end
      default: state <= HALT;              // HALT: stays here until reset
    endcase
  end
end
\end{lstlisting}

Les opérations de l'extension M ont été regroupées pour limiter le matériel. Un seul
multiplieur signé de $33\times33$ bits sert aux quatre multiplications : le bit ajouté
devant chaque opérande vaut son bit de signe ou 0 selon la variante. Un seul diviseur non
signé travaille sur les valeurs absolues, et les signes du quotient et du reste sont
rétablis ensuite. Avec cette méthode, le cas $-2^{31}/(-1)$ donne directement $-2^{31}$ et un
reste nul, comme l'impose la spécification ; seule la division par zéro est traitée à part.

\subsection{Validation par simulation}

Le système \texttt{systeme\_fpga.vhd} relie ce processeur au bloc \texttt{quadram}, réutilisé
sans modification. Les quatre programmes des parties précédentes donnent les mêmes résultats
qu'avec le modèle de la partie~\ref{sec:systeme} (tableau~\ref{tab:fpga-sim}) ; seuls les
instants changent, puisque chaque instruction dure maintenant deux ou trois périodes. La
figure~\ref{fig:p6-systeme} montre la fin de \texttt{test\_memoire} : les états 0, 1 et 2
(\texttt{FETCH}, \texttt{EXEC}, \texttt{MEM}) alternent, l'écriture au mot 1023 est
relue juste après, puis \texttt{ebreak} fait passer le processeur dans l'état 3
(\texttt{HALT}).

\begin{table}[H]
\centering
\begin{tabular}{lrl}
\toprule
Programme & Fin & Résultats \\
\midrule
\texttt{instr2} & 140~ns & x1 = 1, x2 = 6, mot 65 = 6 \\
\texttt{test\_base} & 340~ns & registres du tableau~\ref{tab:testbase}, mots 96 et 104 = $-7$ et 107 \\
\texttt{test\_rv32im} & 380~ns & valeurs du tableau~\ref{tab:rv32imres} \\
\texttt{test\_memoire} & 540~ns & x20 à x23 = 11, 22, 33, 44 ; x24 = 110 ; x25 = 44 ; x26 = 110 \\
\bottomrule
\end{tabular}
\caption{Programmes de test sur le processeur synthétisable (horloge de 10~ns)}
\label{tab:fpga-sim}
\end{table}

\capture{p6-systeme.png}{ModelSim, dossier \texttt{exo/p6}, \texttt{systeme\_fpga\_tb} :
fin de \texttt{test\_memoire}}
{Fin de \texttt{test\_memoire} sur le processeur synthétisable : états de la machine,
accès à \texttt{quadram} et arrêt sur \texttt{ebreak}}{fig:p6-systeme}

\subsection{Implantation sur la carte}

Le niveau supérieur \texttt{top\_fpga.vhd} adapte le système à la carte. Les noms des ports
reprennent ceux du tableau de broches du constructeur.
\begin{itemize}
\item \textbf{Horloge} : le chemin le plus long passe par le diviseur combinatoire de
l'extension M, ce qui limite la fréquence maximale à 5,71~MHz. L'horloge de 50~MHz de la
carte est donc divisée par 16, soit 3,125~MHz. Le fichier \texttt{.sdc} déclare les deux
horloges (\texttt{create\_clock} et \texttt{create\_generated\_clock}).
\item \textbf{Reset} : le bouton RESET, actif à 0, est resynchronisé sur l'horloge du
système ; le reset est aussi actif pendant les deux premières périodes après la
configuration.
\item \textbf{Affichage} : les six afficheurs sept segments, à anode commune, montrent un
registre en décimal (conversion par la méthode \og double dabble\fg{}). Par défaut il s'agit
de x24, la somme des quatre mots relus ; les touches KEY0 à KEY3 affichent x20 à x23,
c'est-à-dire les mots relus dans chacune des quatre RAM.
\item \textbf{LED} : LED0 indique que le programme est terminé, LED1 qu'une instruction
inconnue a été rencontrée, et LED3 clignote tant que l'horloge fonctionne.
\end{itemize}

Un testbench, \texttt{top\_fpga\_tb.vhd}, lit les afficheurs comme on les verrait sur la
carte (figure~\ref{fig:p6-afficheur}) : \og 110\fg{} sans touche, \og 11\fg{} avec KEY0 et
\og 44\fg{} avec KEY3. Le tableau~\ref{tab:quartus-p6} et la figure~\ref{fig:p6-quartus}
résument la compilation dans Quartus :
les quatre \texttt{Sync\_Ram} sont reconnues comme des mémoires internes de
$1024\times32$~bits de quatre blocs M9K chacune, la ROM occupe un bloc, et le bus trois
états de \texttt{quadram} est transformé en logique. Toutes les marges de temps sont
positives.

\capture{p6-afficheur.png}{Transcript ModelSim, dossier \texttt{exo/p6}, \texttt{top\_fpga\_tb}}
{Lecture des afficheurs simulée : x24 sans touche, x20 avec KEY0, x23 avec KEY3}{fig:p6-afficheur}

\begin{table}[H]
\centering
\begin{tabular}{lrr}
\toprule
Ressource & Utilisé & Disponible \\
\midrule
Éléments logiques & 5\,577 (54~\%) & 10\,320 \\
Registres & 1\,050 & 10\,320 \\
Bits de mémoire & 139\,264 (33~\%) & 423\,936 \\
Blocs M9K & 17 & 46 \\
Multiplieurs 9 bits & 8 & 46 \\
Broches & 24 & 180 \\
\midrule
Fréquence maximale (modèle lent, 85~°C) & 5,71~MHz & \\
Fréquence utilisée & 3,125~MHz & \\
\bottomrule
\end{tabular}
\caption{Résultats de compilation Quartus pour l'EP4CE10F17C8 (partie 6)}
\label{tab:quartus-p6}
\end{table}

\capture{p6-quartus.png}{Capture de Quartus II 13.1 : projet \texttt{exo/p6/top\_fpga.qpf}
compilé, rapport \og Flow Summary\fg{}}
{Le projet \texttt{top\_fpga} compilé dans Quartus II 13.1}{fig:p6-quartus}

Le fichier de configuration \texttt{output\_files/top\_fpga.sof} est ensuite chargé dans le
FPGA par le Programmer de Quartus, en mode JTAG, avec le câble USB-Blaster
(figure~\ref{fig:p6-programmer}). Cette configuration est volatile : elle disparaît quand la
carte est mise hors tension.

\capture[0.6\linewidth]{p6-programmer.png}{Capture du Programmer de Quartus : USB-Blaster [USB-0], mode JTAG,
fichier \texttt{top\_fpga.sof}, progression à 100~\% (\og Successful\fg{})}
{Programmation du FPGA avec le Programmer de Quartus}{fig:p6-programmer}

Après programmation par JTAG, la carte affiche 110 et LED0 (marquée LED1 sur la carte)
s'allume : le programme \texttt{test\_memoire} s'est exécuté sur le FPGA, avec les quatre
RAM réalisées dans les blocs M9K. Les touches KEY0 à KEY3 affichent 11, 22, 33 et 44, les
mots relus dans chacune des quatre RAM, et tant que le bouton RESET est enfoncé les registres
sont à 0 (figure~\ref{fig:p6-carte}).

\begin{figure}[H]
\centering
\photo{p6_110.jpg}{sans touche : x24 = 110}\hfill
\photo{p6_11.jpg}{KEY0 : x20 = 11}\hfill
\photo{p6_22.jpg}{KEY1 : x21 = 22}\\[3mm]
\photo{p6_33.jpg}{KEY2 : x22 = 33}\hfill
\photo{p6_44.jpg}{KEY3 : x23 = 44}\hfill
\photo{p6_0.jpg}{RESET enfoncé : 0}
\caption{La carte programmée avec \texttt{test\_memoire}}
\label{fig:p6-carte}
\end{figure}
